\documentclass[11pt]{article}

\usepackage[letterpaper,margin=1in]{geometry}
\usepackage{amsmath,amssymb}
\usepackage{booktabs}
\usepackage{longtable}
\usepackage{array}
\usepackage{xcolor}
\usepackage{listings}
\usepackage{url}
\usepackage{hyperref}
\usepackage{titlesec}
\usepackage{parskip}
\usepackage{fancyhdr}
\usepackage{enumitem}

% Hyperref colors
\hypersetup{
  colorlinks=true,
  linkcolor=black,
  urlcolor=blue!50!black,
  citecolor=black,
  pdftitle={KV Cache Engine --- Interface Specification},
  pdfauthor={LonghornSilicon},
  pdfsubject={ISA Specification, Version kv-isa-0.1}
}

% Headings
\titleformat{\section}{\normalfont\Large\bfseries}{\thesection.}{0.6em}{}
\titleformat{\subsection}{\normalfont\large\bfseries}{\thesubsection}{0.6em}{}
\titleformat{\subsubsection}{\normalfont\normalsize\bfseries}{\thesubsubsection}{0.6em}{}

% Headers / footers
\pagestyle{fancy}
\fancyhf{}
\lhead{\small KV Cache Engine ISA}
\rhead{\small \texttt{kv-isa-0.1}}
\cfoot{\thepage}
\renewcommand{\headrulewidth}{0.4pt}

% Code listings
\definecolor{codebg}{gray}{0.97}
\definecolor{codekey}{rgb}{0.13,0.29,0.53}
\definecolor{codecom}{rgb}{0.25,0.5,0.25}
\lstset{
  basicstyle=\ttfamily\small,
  backgroundcolor=\color{codebg},
  frame=single,
  framerule=0.4pt,
  rulecolor=\color{black!30},
  xleftmargin=0.5em,
  xrightmargin=0.5em,
  aboveskip=0.6em,
  belowskip=0.6em,
  showstringspaces=false,
  breaklines=true,
  keywordstyle=\color{codekey}\bfseries,
  commentstyle=\color{codecom}\itshape,
  numbers=none
}

\setlist{itemsep=1pt,topsep=2pt}

\begin{document}

%==============================================================================
% Title page
%==============================================================================
\thispagestyle{empty}
\begin{center}
\vspace*{1.2in}
{\Huge\bfseries KV Cache Engine}\\[0.4em]
{\LARGE\bfseries Interface Specification}\\[1.2em]

{\large \textsc{LonghornSilicon LLM Inference Accelerator}}\\[0.4em]
{\large Block 2 of 4 --- KV Cache Compression / Decompression}\\[2em]

\rule{0.7\textwidth}{0.4pt}\\[1em]

\begin{tabular}{rl}
\textbf{Version}:    & \texttt{kv-isa-0.1} (first public draft) \\
\textbf{Date}:       & 2026-05-14 \\
\textbf{Status}:     & Pre-tape-out stub for compiler integration \\
\textbf{Author}:     & Chaithu Talasila, The University of Texas at Austin \\
\textbf{Reference}:  & \texttt{LonghornSilicon/kv-cache-engine} \\
\end{tabular}\\[1em]
\rule{0.7\textwidth}{0.4pt}

\vfill
\begin{minipage}{0.85\textwidth}
\small
\textbf{Scope.} This document describes the externally-visible interface to the
\texttt{kv\_cache\_engine} block as it will appear on the chip, on an
FPGA prototype, or in the bit-accurate software model.
A compiler backend targeting the KV cache engine writes against this
interface; the hardware implementation must conform to it. This specification
is stable for the KV cache engine block only and will be unified
with the rest of the LonghornSilicon ISA when the Token
Importance Unit and Memory Hierarchy Controller blocks land.
\end{minipage}

\vspace{0.3in}
\end{center}

\clearpage

%==============================================================================
% TOC
%==============================================================================
\tableofcontents
\clearpage

%==============================================================================
% 1. Block overview
%==============================================================================
\section{Block Overview}
\label{sec:overview}

The KV cache engine manages on-chip storage of key and value tensors for
transformer attention, with lossy compression to reduce DRAM bandwidth.
It implements TurboQuant+ (\texttt{turbo4} configuration):

\begin{itemize}
\item \textbf{Key path}: PolarQuant (3-bit) + QJL (1-bit) = 4.25 bits per value (bpv)
\item \textbf{Value path}: PolarQuant (3-bit) only $\approx$ 3.0 bpv
\end{itemize}

\noindent The compression pipeline consists of:
\begin{enumerate}
\item \textbf{Norm extraction}: L2 norm of the input vector, stored as metadata.
\item \textbf{Rotation}: Walsh-Hadamard Transform with random sign flips
      (deterministic from seed). Self-inverse, $O(n \log n)$ butterfly structure.
\item \textbf{Quantization}: 3-bit optimal Lloyd-Max scalar quantization on the
      rotated coordinates, producing centroid indices.
\item \textbf{QJL projection} (K path only): Quantized Johnson-Lindenstrauss
      transform on the residual. Random $\pm 1$ projection matrix from seed,
      producing 1-bit signs per dimension.
\item \textbf{Packing}: Bit-packing of norm, indices, residual norm, and signs
      into a compact SRAM word.
\end{enumerate}

\noindent Decompression reverses the pipeline: unpack, centroid lookup, QJL
correction (K only), inverse WHT, norm rescaling.

\medskip
\noindent\textbf{Compression ratio:} K = 3.56$\times$, V = 4.92$\times$, combined $\geq$ 3$\times$.\\
\textbf{Footprint:} $\sim$6400 FFs estimated (full pipeline); $\sim$93 FFs current top-level.\\
\textbf{Bit-exact reference:} \path{sw/reference_model/kv_cache_engine_ref.py}
and \path{sw/reference_model/kv_cache_engine_ref.{hpp,cpp}}.

%==============================================================================
% 2. Address space and memory map
%==============================================================================
\section{Address Space and Memory Map}
\label{sec:memmap}

The block exposes a 256-byte AXI-Lite slave register window. All registers
are 32-bit, word-aligned. The base address is set at chip integration time.

\begin{table}[h]
\centering
\small
\caption{AXI-Lite register window (offsets relative to base address).}
\label{tab:regs}
\begin{tabular}{@{}l l c l p{6cm}@{}}
\toprule
Offset & Name & Access & Reset & Purpose \\
\midrule
\texttt{0x00} & \texttt{CTRL}              & RW    & 0x0   & bit[0]: \texttt{soft\_reset} (write-1 to pulse); bit[1]: \texttt{enable} \\
\texttt{0x04} & \texttt{STATUS}            & R     & 0x1   & bit[0]: \texttt{idle}; bit[1]: reserved; bit[2]: reserved; bit[3]: \texttt{sram\_full} \\
\texttt{0x08} & \texttt{INFO\_DIM}         & R     & (syn) & Synthesis-time \texttt{VECTOR\_DIM} \\
\texttt{0x0C} & \texttt{INFO\_PQ\_BITS}    & R     & (syn) & PolarQuant quantization bits (3) \\
\texttt{0x10} & \texttt{INFO\_QJL\_BITS}   & R     & (syn) & QJL projection bits (1) \\
\texttt{0x14} & \texttt{INFO\_SRAM\_DEPTH} & R     & (syn) & Number of SRAM entries \\
\texttt{0x18} & \texttt{INFO\_CR\_K}       & R     & (syn) & Key compression ratio, fixed-point 8.8 \\
\texttt{0x1C} & \texttt{INFO\_CR\_V}       & R     & (syn) & Value compression ratio, fixed-point 8.8 \\
\texttt{0x20} & \texttt{INFO\_VERSION}     & R     & 0x00010000 & bits[31:24] major, bits[23:16] minor, bits[15:8] patch, bits[7:0] build \\
\texttt{0x24} & \texttt{OCCUPANCY}         & R     & 0x0   & Number of valid entries in SRAM \\
\texttt{0x28} & \texttt{WRITE\_ADDR}       & RW    & 0x0   & Target SRAM address for next write \\
\texttt{0x2C} & \texttt{READ\_ADDR}        & RW    & 0x0   & Target SRAM address for next read \\
\texttt{0x30} & \texttt{KV\_SELECT}        & RW    & 0x0   & bit[0]: 0 = key path, 1 = value path \\
\texttt{0x34} & \texttt{IRQ\_MASK}         & RW    & 0x0   & bits[3:0]: interrupt enable mask \\
\texttt{0x38} & \texttt{IRQ\_STATUS}       & R/W1C & 0x0   & bits[3:0]: interrupt pending; write-1-to-clear \\
\bottomrule
\end{tabular}
\end{table}

\medskip
\noindent\textbf{Conventions.} \texttt{RW} = read/write; \texttt{R} = read-only;
\texttt{W1C} = write-1-to-clear; (syn) = value fixed at synthesis time.
Writes to read-only registers are silently dropped. Reading reserved offsets
returns \texttt{0xDEADBEEF}.

%==============================================================================
% 3. Streaming data interfaces
%==============================================================================
\section{Streaming Data Interfaces}
\label{sec:stream}

Two AXI-Stream interfaces carry the vector data. The AXI-Lite window in
\S\ref{sec:memmap} is for control only.

\subsection{\texttt{s\_axis\_kv} --- Vector Write (Slave)}
\label{sec:s_axis_kv}

\begin{table}[h]
\centering
\small
\begin{tabular}{@{}l c c p{8cm}@{}}
\toprule
Signal      & Width                  & Direction & Purpose \\
\midrule
\texttt{tdata}  & \texttt{COORD\_WIDTH} & in        & One signed two's-complement coordinate \\
\texttt{tlast}  & 1                     & in        & Assert on the final coordinate of the vector \\
\texttt{tvalid} & 1                     & in        & Handshake: data is valid \\
\texttt{tready} & 1                     & out       & Handshake: block is ready to accept \\
\texttt{tuser}  & 1                     & in        & 0 = key vector, 1 = value vector \\
\bottomrule
\end{tabular}
\end{table}

\noindent\textbf{Protocol:}
\begin{itemize}
\item A complete vector is exactly \texttt{VECTOR\_DIM} beats. \texttt{tlast}
  asserts on the final beat.
\item The \texttt{tuser} field selects the compression path: key vectors receive
  PolarQuant + QJL compression; value vectors receive PolarQuant only.
\item The target SRAM address must be set via \texttt{WRITE\_ADDR} before
  streaming begins.
\item Backpressure: when the block is not ready (e.g., during compression or
  SRAM write), \texttt{tready} deasserts. Upstream must hold signals stable.
\end{itemize}

\subsection{\texttt{m\_axis\_kv} --- Vector Read (Master)}
\label{sec:m_axis_kv}

\begin{table}[h]
\centering
\small
\begin{tabular}{@{}l c c p{8cm}@{}}
\toprule
Signal      & Width                    & Direction & Purpose \\
\midrule
\texttt{tdata}  & \texttt{COORD\_WIDTH}   & out       & One signed two's-complement decompressed coordinate \\
\texttt{tlast}  & 1                       & out       & Asserted on the final coordinate \\
\texttt{tvalid} & 1                       & out       & Handshake: data is valid \\
\texttt{tready} & 1                       & in        & Handshake: downstream is ready to accept \\
\bottomrule
\end{tabular}
\end{table}

\noindent\textbf{Protocol.} One beat per coordinate, \texttt{VECTOR\_DIM} beats
per vector. Set \texttt{READ\_ADDR} and \texttt{KV\_SELECT} before initiating
a read. The decompression path is selected by \texttt{KV\_SELECT}.

\subsection{Eviction Interface}
\label{sec:evict}

\begin{table}[h]
\centering
\small
\begin{tabular}{@{}l c c p{8cm}@{}}
\toprule
Signal             & Width                         & Direction & Purpose \\
\midrule
\texttt{evict\_needed} & 1                         & out       & Asserted when SRAM is full \\
\texttt{evict\_addr}   & $\lceil\log_2 D\rceil$    & out       & Address of entry to evict \\
\bottomrule
\end{tabular}
\end{table}

\noindent This interface connects to the Memory Hierarchy Controller (block~4).
When \texttt{evict\_needed} asserts, the controller must either spill the entry
at \texttt{evict\_addr} to DRAM or invalidate it before the next write.

%==============================================================================
% 4. Logical operations (compiler-facing)
%==============================================================================
\section{Logical Operations}
\label{sec:ops}

Below are the abstract operations a compiler emits. Each maps to a short
sequence of register writes and stream beats; the C API in \S\ref{sec:capi}
gives the concrete library calls.

\begin{table}[h]
\centering
\small
\begin{tabular}{@{}l l l p{5.5cm}@{}}
\toprule
Operation                     & Inputs             & Outputs          & Description \\
\midrule
\texttt{KV\_QUERY}            & ---                & INFO struct      & Read all \texttt{INFO\_*} registers \\
\texttt{KV\_RESET}            & ---                & ---              & Soft reset: clear SRAM, reset FSM \\
\texttt{KV\_ENABLE}           & ---                & ---              & Set bit[1] of \texttt{CTRL} \\
\texttt{KV\_COMPRESS\_KEY}    & addr, $D$ coords   & ---              & Compress and store a key vector at \texttt{addr} \\
\texttt{KV\_COMPRESS\_VALUE}  & addr, $D$ coords   & ---              & Compress and store a value vector at \texttt{addr} \\
\texttt{KV\_DECOMPRESS\_KEY}  & addr               & $D$ coords       & Read and decompress a key vector from \texttt{addr} \\
\texttt{KV\_DECOMPRESS\_VALUE}& addr               & $D$ coords       & Read and decompress a value vector from \texttt{addr} \\
\texttt{KV\_STATUS}           & ---                & STATUS bits      & Read the \texttt{STATUS} register \\
\texttt{KV\_OCCUPANCY}        & ---                & count            & Read the \texttt{OCCUPANCY} register \\
\bottomrule
\end{tabular}
\end{table}

\subsection{Worked Lowering Example}
\label{sec:lowering}

This section walks through how a compiler lowers a transformer attention
operation into KV cache engine operations.

Suppose the compiler's IR contains:
\begin{lstlisting}[basicstyle=\ttfamily\small\itshape]
%out = transformer_block(%q, %k, %v, %kv_cache)
    : tensor<S x D x f16>, tensor<S x D x f16>, tensor<S x D x f16>,
      kv_cache_handle -> tensor<S x D x f16>
\end{lstlisting}

\textbf{Stage~1: setup.} Query hardware configuration once per compilation:

\begin{lstlisting}[basicstyle=\ttfamily\small]
info = KV_QUERY()       // -> dim=64, pq_bits=3, qjl_bits=1, sram_depth=16
KV_RESET()
KV_ENABLE()
\end{lstlisting}

\textbf{Stage~2: store phase.} For each new token, compress and store its K/V:

\begin{lstlisting}[basicstyle=\ttfamily\small]
for token_idx in new_tokens:
    addr = token_idx % info.sram_depth
    KV_COMPRESS_KEY(addr, k_vectors[token_idx])
    KV_COMPRESS_VALUE(addr, v_vectors[token_idx])
\end{lstlisting}

\textbf{Stage~3: attention phase.} Decompress cached K/V for attention:

\begin{lstlisting}[basicstyle=\ttfamily\small]
for cached_idx in range(num_cached):
    k_hat = KV_DECOMPRESS_KEY(cached_idx)
    v_hat = KV_DECOMPRESS_VALUE(cached_idx)
    score = dot(query, k_hat) / sqrt(D)

    // Route through precision controller (block 1)
    decision = PC_PUSH_TILE(score)
    if decision == FP16:
        out += fp16_matmul(softmax(score), v_hat)
    else:
        out += int8_matmul(softmax(score), v_hat)
\end{lstlisting}

\noindent Note the interaction between block~1 (precision controller) and
block~2 (KV cache engine): the decompressed K/V vectors feed into the
attention score computation, which the precision controller routes.

\subsection{Compiler Binding Patterns}
\label{sec:bindings}

\noindent\textbf{MLIR dialect.} Define \texttt{lhsi.kv.*} ops mirroring the
table in \S\ref{sec:ops}. The lowering pass rewrites \texttt{linalg.attention}
or vendor-equivalent ops to emit \texttt{lhsi.kv.compress\_key},
\texttt{lhsi.kv.decompress\_key}, etc.

\noindent\textbf{TVM Relax / TIR.} Register backend target \texttt{"lhsi-kv"},
pattern-match KV cache operations in the attention sub-graph, emit calls to
TVM-runtime extern functions named for each \texttt{KV\_*} operation.

\noindent\textbf{ONNX.} Add \texttt{CustomOpDomain} with KV cache ops, graph
transform inserts compression/decompression around attention.

\noindent\textbf{Custom IR.} Replace KV cache codegen with \texttt{KV\_*}
operation sequences targeting the C API.

%==============================================================================
% 5. C API stub
%==============================================================================
\section{C API Stub}
\label{sec:capi}

\begin{lstlisting}[language=C]
/* lhsi_kv_cache_engine.h --- LonghornSilicon KV Cache Engine driver API. */

#include <stdint.h>
#include <stdbool.h>
#include <stddef.h>

typedef struct {
    uint32_t vector_dim;
    uint32_t pq_bits;
    uint32_t qjl_bits;
    uint32_t sram_depth;
    uint32_t cr_k_fixed;       /* compression ratio K, fixed-point 8.8 */
    uint32_t cr_v_fixed;       /* compression ratio V, fixed-point 8.8 */
    uint32_t version;
} lhsi_kv_info_t;

typedef struct lhsi_kv_handle lhsi_kv_handle_t;

/* Open / close */
int  lhsi_kv_open (const char *device, lhsi_kv_handle_t **out);
void lhsi_kv_close(lhsi_kv_handle_t *h);

/* Configuration queries */
int  lhsi_kv_query(lhsi_kv_handle_t *h, lhsi_kv_info_t *out);

/* Control */
int  lhsi_kv_reset (lhsi_kv_handle_t *h);
int  lhsi_kv_enable(lhsi_kv_handle_t *h, bool enable);

/* Compress and store a key vector at the given SRAM address. */
int  lhsi_kv_compress_key(lhsi_kv_handle_t *h, uint32_t addr,
                          const int16_t *coords, size_t dim);

/* Compress and store a value vector at the given SRAM address. */
int  lhsi_kv_compress_value(lhsi_kv_handle_t *h, uint32_t addr,
                            const int16_t *coords, size_t dim);

/* Read and decompress a key vector from the given SRAM address. */
int  lhsi_kv_decompress_key(lhsi_kv_handle_t *h, uint32_t addr,
                            int16_t *coords, size_t dim);

/* Read and decompress a value vector from the given SRAM address. */
int  lhsi_kv_decompress_value(lhsi_kv_handle_t *h, uint32_t addr,
                              int16_t *coords, size_t dim);

/* Diagnostics. */
typedef struct {
    bool     idle;
    bool     sram_full;
    uint32_t occupancy;
} lhsi_kv_status_t;
int  lhsi_kv_status(lhsi_kv_handle_t *h, lhsi_kv_status_t *out);
\end{lstlisting}

\noindent
Return codes follow the standard Linux convention: \texttt{0} on success,
\texttt{-errno} on failure.

%==============================================================================
% 6. Compression algorithm
%==============================================================================
\section{Compression Algorithm: TurboQuant+}
\label{sec:algorithm}

TurboQuant+ is a two-stage vector quantization scheme designed for KV cache
compression in transformer inference. The \texttt{turbo4} configuration
provides 4.25 bpv for keys and $\sim$3.0 bpv for values.

\subsection{PolarQuant (Both K and V)}

\begin{enumerate}
\item \textbf{Norm extraction}: Compute $\|x\|_2$ (L2 norm) and normalize:
      $\hat{x} = x / \|x\|_2$.
\item \textbf{Rotation}: Apply random sign flips $s_i \in \{-1, +1\}$ (from
      LCG seed), then Walsh-Hadamard Transform:
      $y = \mathrm{WHT}(\hat{x} \odot s)$.
      This randomizes the coordinate distribution toward i.i.d.\ Gaussian.
\item \textbf{Quantization}: 3-bit optimal Lloyd-Max scalar quantization of
      each rotated coordinate $y_i$, using 8 centroids symmetric around zero:
      $\pm 0.2451\sigma$, $\pm 0.7560\sigma$, $\pm 1.3439\sigma$,
      $\pm 2.1520\sigma$ where $\sigma = 1/\sqrt{D}$.
\end{enumerate}

\subsection{QJL Correction (K Path Only)}

After PolarQuant, the quantization residual $r = y - \hat{y}$ is projected
through a random $\pm 1$ matrix $R \in \{-1,+1\}^{D \times D}$ (from
seed + \texttt{0xDEADBEEF}):
\[
z_j = \mathrm{sign}\bigl(\textstyle\sum_{i=1}^{D} R_{ji} \cdot r_i\bigr)
\]
The 1-bit signs $z_j$ preserve inner products in expectation (Johnson-Lindenstrauss
property). During decompression, the QJL correction is added back to the
dequantized PolarQuant output, improving attention score accuracy for keys.

\subsection{Bit Layout}

\begin{table}[h]
\centering
\small
\begin{tabular}{@{}l l l l@{}}
\toprule
Path & Fields & Bits & Compression \\
\midrule
Key  & norm(16) + indices($D \times 3$) + res\_norm(16) + signs($D \times 1$) & $32 + 4D$ & 3.56$\times$ \\
Value & norm(16) + indices($D \times 3$) & $16 + 3D$ & 4.92$\times$ \\
\bottomrule
\end{tabular}
\end{table}

\noindent For $D = 64$: key = 288 bits, value = 208 bits,
vs.\ uncompressed $D \times 16 = 1024$ bits.

%==============================================================================
% 7. Synthesis-time configuration
%==============================================================================
\section{Synthesis-Time Configuration}
\label{sec:syntime}

\begin{table}[h]
\centering
\small
\begin{tabular}{@{}l c l p{5cm}@{}}
\toprule
Parameter             & Default & Range                   & Notes \\
\midrule
\texttt{VECTOR\_DIM}    & 64  & 16, 32, 64, 128          & Must be a power of two (WHT constraint) \\
\texttt{PQ\_BITS}       & 3   & 2, 3, 4                  & Lloyd-Max quantization bits \\
\texttt{QJL\_BITS}      & 1   & 1                        & Fixed at 1 in \texttt{kv-isa-0.1} \\
\texttt{SRAM\_DEPTH}    & 16  & 4--4096                  & Behavioral reg array for Sky130; real macros at 16FFC \\
\texttt{COORD\_WIDTH}   & 16  & 8, 12, 16                & Fixed-point coordinate width \\
\texttt{COORD\_FRAC}    & 12  & 4--14                    & Fractional bits in coordinate representation \\
\texttt{NORM\_WIDTH}    & 16  & 8, 12, 16                & Norm storage width \\
\texttt{ROTATION\_SEED} & 42  & any 32-bit               & Deterministic LCG seed for sign flips and QJL matrix \\
\bottomrule
\end{tabular}
\end{table}

%==============================================================================
% 8. Bit-Accurate Reference
%==============================================================================
\section{Bit-Accurate Reference}
\label{sec:ref}

Three implementations are provided, all bit-exact with each other:

\begin{enumerate}
\item \textbf{Python}: \path{sw/reference_model/kv_cache_engine_ref.py} ---
      golden reference with fixed-point arithmetic.
\item \textbf{C++}: \path{sw/reference_model/kv_cache_engine_ref.{hpp,cpp}} ---
      same fixed-point math, with \texttt{extern "C"} API using opaque handles.
\item \textbf{SystemVerilog}: \path{rtl/kv_cache_engine.sv} + sub-modules ---
      synthesizable RTL.
\end{enumerate}

\noindent Usage:
\begin{lstlisting}[language=Python]
from kv_cache_engine_ref import KVCacheEngine, KVCacheEngineInfo

engine = KVCacheEngine()
ck = engine.compress_key(vector)
reconstructed = engine.decompress_key(ck)
\end{lstlisting}

%==============================================================================
% 9. Integration phases
%==============================================================================
\section{Integration Phases}
\label{sec:phases}

\begin{table}[ht]
\centering
\small
\renewcommand{\arraystretch}{1.2}
\begin{tabular}{@{}c p{3.0cm} p{4.0cm} p{5.0cm}@{}}
\toprule
Phase & Timeline & Compiler targets & We provide \\
\midrule
0 & now & Bit-accurate Python and C++ reference models & Models + ISA + tests \\
1 & when FPGA arrives & AXI-Lite + AXI-Stream on Zynq UltraScale+ & Vivado bitstream + driver \\
2 & after all 4 blocks complete & Multi-block FPGA project & Integrated attention pipeline \\
3 & post-tape-out (2027+) & PCIe-attached custom chip & TSMC 16FFC silicon + Linux driver \\
\bottomrule
\end{tabular}
\end{table}

%==============================================================================
% 10. Interaction with other blocks
%==============================================================================
\section{Interaction with Other Blocks}
\label{sec:interaction}

\begin{itemize}
\item \textbf{Block 1 (Precision Controller)}: Receives decompressed K/V vectors
      from this block for attention score computation. The precision controller
      decides INT8 vs FP16 routing based on the decompressed scores.
\item \textbf{Block 3 (Token Importance Unit)}: Provides token importance scores
      that inform eviction decisions when the SRAM is full.
\item \textbf{Block 4 (Memory Hierarchy Controller)}: Receives
      \texttt{evict\_needed} / \texttt{evict\_addr} signals when SRAM is full.
      Manages spilling compressed entries to DRAM and fetching them back.
\end{itemize}

%==============================================================================
% 11. Open questions
%==============================================================================
\section{Open Questions}
\label{sec:open}

\begin{enumerate}
\item Should the compression path support variable-rate quantization (2/3/4-bit
      PolarQuant selected per-vector based on norm)?
\item Should the QJL projection dimension be configurable (currently $D \times D$,
      could be $D \times D'$ with $D' < D$ for reduced storage)?
\item Should the block support direct SRAM-to-SRAM copy for cache migration?
\item Are the current fixed-point widths sufficient for larger models
      ($D = 128, 256$)?
\item Should the eviction policy be configurable via registers (LRU, FIFO,
      importance-based)?
\end{enumerate}

%==============================================================================
% 12. Change log
%==============================================================================
\section{Change Log}
\label{sec:changelog}

\begin{itemize}
\item \texttt{kv-isa-0.1} (2026-05-14): First public draft. Stable for the
      KV cache engine block only.
\end{itemize}

\vspace{0.5in}
\noindent\rule{\textwidth}{0.4pt}\\
\noindent\small\textit{LonghornSilicon KV Cache Engine --- Interface Specification.
This document is open and may be redistributed. The TSMC 16FFC PDK referenced
in the integration phases is NDA-protected and is not part of this document
or the public repository.}

\end{document}
